% =====================================================================
%  Chapitre 3 — Sprint 1 : Authentification, multi-entreprises
%                          et gestion des utilisateurs
% =====================================================================
\chapter{Sprint~1 : Authentification, multi-entreprises et gestion des utilisateurs}
\label{chap:sprint1}

\section*{Introduction}
Ce premier sprint pose les fondations de la plateforme : la manière dont un
utilisateur accède au système, dont ses données sont cloisonnées par
entreprise, et dont les membres d'une même organisation sont gérés. Nous
présentons successivement les objectifs du sprint, le raffinement des cas
d'utilisation concernés, les diagrammes de séquence associés, puis la
réalisation obtenue.

\section{Objectifs du sprint}
Ce sprint couvre les epics~1 à~3 du backlog produit. Le
tableau~\ref{tab:backlog-s1} en détaille les user stories.

\begin{table}[H]
\centering
\caption{Backlog du sprint~1}
\label{tab:backlog-s1}
\small
\begin{tabularx}{\linewidth}{clX}
\toprule
\textbf{ID} & \textbf{User story} & \textbf{Description} \\
\midrule
1.1 & S'inscrire & En tant qu'utilisateur, je veux créer un compte et une entreprise afin de disposer d'un espace de travail. \\
1.2 & Se connecter & En tant qu'utilisateur, je veux m'authentifier de façon sécurisée pour accéder à la plateforme. \\
1.3 & Gérer mon profil & En tant qu'utilisateur, je veux consulter et modifier mes informations personnelles. \\
2.1 & Isolation des données & En tant qu'entreprise, je veux que mes modèles et mes membres soient invisibles pour les autres entreprises. \\
3.1 & Inviter un membre & En tant qu'administrateur, je veux inviter des collaborateurs à rejoindre mon entreprise. \\
3.2 & Rejoindre une entreprise & En tant qu'invité, je veux accepter une invitation pour intégrer une entreprise. \\
3.3 & Gérer l'équipe & En tant qu'administrateur, je veux lister les membres et leur attribuer un rôle (administrateur, éditeur, marketing). \\
3.4 & Superviser les entreprises & En tant que super-administrateur, je veux consulter l'ensemble des entreprises et de leurs abonnements. \\
\bottomrule
\end{tabularx}
\end{table}

\section{Raffinement et description des cas d'utilisation}
Nous précisons dans cette section le fonctionnement des services
implémentés, en détaillant les processus clés et les interactions entre
acteurs et système.

\subsection{Cas d'utilisation « Authentification et inscription »}
La figure~\ref{fig:uc-auth} illustre les cas d'utilisation relatifs à
l'inscription et à l'authentification. L'inscription crée conjointement un
compte utilisateur et l'entreprise associée~; la connexion s'effectue par
identifiants et délivre un jeton de session.

% [FIGURE : diagramme de cas d'utilisation Authentification / Inscription]
\figplaceholder{Cas d'utilisation « Authentification et inscription »}{fig:uc-auth}{7cm}

Le tableau~\ref{tab:uc-inscription} décrit textuellement le scénario
« S'inscrire ».

\begin{table}[H]
\centering
\caption{Description textuelle du cas d'utilisation « S'inscrire »}
\label{tab:uc-inscription}
\small
\begin{tabularx}{\linewidth}{lX}
\toprule
\textbf{Cas d'utilisation} & S'inscrire \\
\textbf{Acteur} & Utilisateur \\
\textbf{Préconditions} & L'utilisateur n'est pas connecté. \\
\textbf{Postconditions} & Un compte et une entreprise sont créés~; l'utilisateur est authentifié. \\
\midrule
\textbf{Scénario nominal} &
1.~L'utilisateur accède au formulaire d'inscription.\newline
2.~Il saisit ses informations et le nom de son entreprise.\newline
3.~Le système valide les données fournies.\newline
4.~Le système crée l'entreprise puis le compte associé.\newline
5.~Le système génère un jeton de session et redirige l'utilisateur vers le tableau de bord. \\
\midrule
\textbf{Scénario alternatif} &
3.a.~Si l'adresse e-mail ou le nom d'entreprise existe déjà, le système en informe l'utilisateur et l'invite à corriger sa saisie. \\
\bottomrule
\end{tabularx}
\end{table}

\subsection{Cas d'utilisation « Gestion des utilisateurs »}
La figure~\ref{fig:uc-users} présente les cas d'utilisation relatifs à la
gestion de l'équipe : invitation d'un membre, acceptation d'une invitation,
consultation des membres et attribution des rôles.

% [FIGURE : diagramme de cas d'utilisation Gestion des utilisateurs]
\figplaceholder{Cas d'utilisation « Gestion des utilisateurs »}{fig:uc-users}{7cm}

Le tableau~\ref{tab:uc-invite} décrit le scénario « Inviter un membre ».

\begin{table}[H]
\centering
\caption{Description textuelle du cas d'utilisation « Inviter un membre »}
\label{tab:uc-invite}
\small
\begin{tabularx}{\linewidth}{lX}
\toprule
\textbf{Cas d'utilisation} & Inviter un membre \\
\textbf{Acteur} & Administrateur d'entreprise \\
\textbf{Préconditions} & L'administrateur est connecté. \\
\textbf{Postconditions} & Une invitation est émise et le futur membre est notifié. \\
\midrule
\textbf{Scénario nominal} &
1.~L'administrateur ouvre la page de gestion de l'équipe.\newline
2.~Il saisit l'adresse e-mail de l'invité et sélectionne un rôle.\newline
3.~Le système crée une invitation rattachée à l'entreprise et envoie un lien à l'invité.\newline
4.~L'invité suit le lien, définit son mot de passe et rejoint l'entreprise. \\
\midrule
\textbf{Scénario alternatif} &
2.a.~Si l'adresse correspond déjà à un membre de l'entreprise, le système signale que l'utilisateur en fait partie. \\
\bottomrule
\end{tabularx}
\end{table}

\section{Diagrammes de séquence}

\subsection{Séquence « S'inscrire »}
La figure~\ref{fig:seq-inscription} décrit les échanges entre le client, la
passerelle API et le service d'authentification lors d'une inscription. La
passerelle traduit la requête REST en appel gRPC~; le service valide les
données, crée l'entreprise et le compte, puis renvoie un jeton porté par un
\emph{cookie} sécurisé.

% [FIGURE : diagramme de séquence « S'inscrire »]
\figplaceholder{Diagramme de séquence « S'inscrire »}{fig:seq-inscription}{7cm}

\subsection{Séquence « Inviter un membre »}
La figure~\ref{fig:seq-invite} illustre le flux d'invitation d'un membre,
depuis la demande de l'administrateur jusqu'à l'acceptation par l'invité.

% [FIGURE : diagramme de séquence « Inviter un membre »]
\figplaceholder{Diagramme de séquence « Inviter un membre »}{fig:seq-invite}{7cm}

\section{Réalisation}
Cette section présente les résultats obtenus au cours du sprint, tant sur le
plan technique que du point de vue des interfaces.

\subsection{Service d'authentification et modèle multi-entreprises}
Le service d'authentification a été développé avec NestJS selon le patron
CQRS et expose ses opérations en gRPC. Il prend en charge l'inscription, la
connexion, la validation de jeton, l'invitation et l'acceptation de membres,
la consultation et la mise à jour du profil, ainsi que le listage des
membres d'une entreprise. Les données sont persistées dans la base
\texttt{auth\_db} via l'ORM TypeORM.

Le c\oe{}ur du modèle multi-entreprises repose sur la notion de
\emph{tenant} (entreprise). Chaque compte est rattaché à un
\texttt{tenant\_id}, qui est embarqué dans la charge utile du jeton JWT aux
côtés de l'identifiant utilisateur, de l'adresse e-mail et du rôle. Ce
\texttt{tenant\_id} est systématiquement propagé aux services applicatifs
et sert de filtre d'isolation : aucune requête ne peut accéder aux données
d'une autre entreprise. La passerelle API constitue l'unique surface REST
publique et vérifie le jeton, transporté par un \emph{cookie}, avant de
relayer chaque appel en gRPC vers le service concerné.

\subsection{Inscription et connexion}
Les figures~\ref{fig:ui-register} et~\ref{fig:ui-login} présentent
respectivement les interfaces d'inscription et de connexion. L'inscription
recueille les informations de l'utilisateur ainsi que le nom de son
entreprise, puis crée l'espace de travail correspondant. La connexion
authentifie l'utilisateur et ouvre une session sécurisée.

% [FIGURE : capture — interface d'inscription]
\figplaceholder{Interface d'inscription}{fig:ui-register}{6cm}

% [FIGURE : capture — interface de connexion]
\figplaceholder{Interface de connexion}{fig:ui-login}{6cm}

\subsection{Gestion du profil}
L'interface de profil (figure~\ref{fig:ui-profile}) permet à l'utilisateur
de consulter et de modifier ses informations personnelles.

% [FIGURE : capture — interface de profil]
\figplaceholder{Interface de gestion du profil}{fig:ui-profile}{6cm}

\subsection{Gestion de l'équipe et des rôles}
L'interface de gestion de l'équipe (figure~\ref{fig:ui-team}) permet à
l'administrateur d'inviter de nouveaux membres, de consulter la liste des
collaborateurs et de leur attribuer un rôle. Trois rôles ont été définis~:
\textbf{administrateur} (droits étendus sur l'entreprise),
\textbf{éditeur} (création et édition de modèles) et \textbf{marketing}
(éditeur habilité à maintenir la galerie de modèles prédéfinis de
l'entreprise). Un rôle machine à machine complète ce dispositif pour les
clients applicatifs authentifiés par clé d'API.

% [FIGURE : captures — interfaces d'invitation et de gestion de l'équipe]
\figplaceholder{Interfaces d'invitation et de gestion de l'équipe}{fig:ui-team}{6cm}

\subsection{Espace de super-administration}
Un espace de super-administration, réservé à l'éditeur de la plateforme,
offre une vue transverse sur l'ensemble des entreprises et de leurs
abonnements. Il permet notamment d'attribuer un plan à une entreprise. Cet
espace, dont les fonctions liées à la facturation sont approfondies au
chapitre~\ref{chap:sprint4}, est présenté à la figure~\ref{fig:ui-superadmin}.

% [FIGURE : capture — tableau de bord super-administration]
\figplaceholder{Tableau de bord de super-administration}{fig:ui-superadmin}{6cm}

\section*{Conclusion}
Ce sprint a établi le socle de la plateforme : un service d'authentification
robuste, un modèle multi-entreprises garantissant l'isolation des données
par \texttt{tenant\_id}, et une gestion complète des utilisateurs et des
rôles. Ces fondations conditionnent l'ensemble des fonctionnalités
ultérieures. Le sprint suivant s'appuie sur cet acquis pour construire le
c\oe{}ur métier de la plateforme : l'éditeur de modèles multi-canal.
